\documentclass[journal,compsoc]{IEEEtran}\newcommand{\journal}{true}
%\documentclass[conference]{IEEEtran}\newcommand{\journal}{false}
\ifCLASSOPTIONcompsoc
  % IEEE Computer Society needs nocompress option
  % requires cite.sty v4.0 or later (November 2003)
  \usepackage[nocompress]{cite}
\else
  % normal IEEE
   \usepackage{cite}
\fi
%\providecommand{\tabularnewline}{\\}

\usepackage{indentfirst}

%%--------------------Colocados por minha vontade própria-----------------------
\usepackage{marvosym}
\usepackage[utf8x]{inputenc}
\usepackage[english,brazil]{babel}	
\usepackage[T1]{fontenc}
\usepackage{amsmath}
\usepackage{wasysym}				% setas
\usepackage{picinpar}				% figuras onde eu quero, problemas com referencias
\usepackage{sidecap}				% legenda ao lado da figura
\usepackage[small,bf]{caption}

\usepackage{multicol}
\usepackage{tikz}
\usepackage{circuitikz}
\usepackage{siunitx}
\usepackage{lipsum}
\usepackage{csquotes}
\usepackage{alltt}

\usepackage{ifsym}
\usepackage{pifont}
\usepackage{manfnt}
\usepackage{stmaryrd}
\usepackage{ulsy}

%\usepackage{fancyhdr}
%\lhead{O que quero no cabeçalho parte esquerda}
%\chead{O que quero no cabeçalho parte central}
%\rhead{O que quero no cabeçalho parte direita}
%\lfoot{O que quero no rodapé parte esquerda}
%\cfoot{O que quero no rodapé parte central}
%\rfoot{O que quero no rodapé parte direita}

\input nomes




%%------------------------------------------------------------------------------
%%	Resolveu o problema de a linguagem não estar declarada
\makeatletter
\def\markboth#1#2{\def\leftmark{\@IEEEcompsoconly{\sffamily}\MakeUppercase{\protect#1}}%
\def\rightmark{\@IEEEcompsoconly{\sffamily}\MakeUppercase{\protect#2}}}
\makeatother
%%------------------------------------------------------------------------------
%%----------------Para inclusão de algoritmos em C/C++/Assembly---------------------------
\usepackage{listings}
%\lstnewenvironment{C}
\lstset{%
language=C,							%linguagem
numbers=left,						%posição dos números
stepnumber=1,						%frequencia de aparição dos números
numbersep=3pt,
basewidth={0.6em,0.45em},
%fontadjust=true,
mathescape=true,
tabsize=4,
commentstyle=\color{blue},
literate={á}{{\'a}}1 {ã}{{\~a}}1 {é}{{\'e}}1 {õ}{{\~o}}1 {í}{{\'i}}1 {ó}{{\'o}}1 {ú}{{\'u}}1 {ç}{{\c c}}1,
breaklines=true,
showstringspaces=false,
stringstyle=\color{cyan},
basicstyle=\small\ttfamily}
% Inclua o source com: \lstinputlisting[language=C]{src/src.c}

\lstdefinelanguage{MSP430}
{ %% Se estiver faltando alguma palavra reservada, insira ela aqui
morekeywords={RRC,SWPB,RRA,SXT,PUSH,CALL,RETI,JNE,JNZ,JEQ,
Z,JNC,JLO,JC,JHS,JN,JGE,JL,JMP,NOP,POP,BR,
RET,CLRC,SETC,CLRZ,SETZ,CLRN,SETN,DINT,EINT,
RLA,RLC,ADD,ADDC,INV,XOR,CLR,MOV,TST,CMP,DEC,
SUB,SUBC,DECD,INC,INCD,SBC,MOV,DADD,BIT,
BIC,BIS,XOR,AND,MOV.B,MOV.W},
sensitive=false,
morecomment=[l]{//},
morecomment=[l]{;},
morecomment=[s]{/*}{*/},
}
% Inclua o source com: \lstinputlisting[language=MSP430]{src/src.asm}

\lstdefinestyle{Bash}
{language=bash,
keywordstyle=\color{blue},
basicstyle=\small\ttfamily,
morekeywords={user@DESKTOP},
alsoletter={:~$},
morekeywords=[2]{peter@kbpet:},
keywordstyle=[2]{\color{red}},
literate={\$}{{\textcolor{red}{\$}}}1 
         {:}{{\textcolor{red}{:}}}1
         {~}{{\textcolor{red}{\textasciitilde}}}1,
}

%
%%------------------------------------------------------------------------------
%%---------------------Includes importantes para suporte------------------------
\input link							% Para verificar links de citações e outras configurações
%%------------------------------------------------------------------------------
%%------------------------------------------------------------------------------



% *** GRAPHICS RELATED PACKAGES ***
%

\ifCLASSINFOpdf
\usepackage{tikz}
\usetikzlibrary{shapes,arrows}
\usetikzlibrary{calc,decorations.pathmorphing,patterns}
\pgfdeclaredecoration{penciline}{initial}{
    \state{initial}[width=+\pgfdecoratedinputsegmentremainingdistance,
    auto corner on length=1mm,]{
        \pgfpathcurveto%
        {% From
            \pgfqpoint{\pgfdecoratedinputsegmentremainingdistance}
                      {\pgfdecorationsegmentamplitude}
        }
        {%  Control 1
        \pgfmathrand
        \pgfpointadd{\pgfqpoint{\pgfdecoratedinputsegmentremainingdistance}{0pt}}
                    {\pgfqpoint{-\pgfdecorationsegmentaspect
                     \pgfdecoratedinputsegmentremainingdistance}%
                               {\pgfmathresult\pgfdecorationsegmentamplitude}
                    }
        }
        {%TO 
        \pgfpointadd{\pgfpointdecoratedinputsegmentlast}{\pgfpoint{1pt}{1pt}}
        }
    }
    \state{final}{}
}


% \usepackage{tikzgraphicx}

\tikzstyle{cloud} = [draw, ellipse,fill=blue!20, node distance=3cm,
    minimum height=2em]
\tikzstyle{estado} = [draw, rectangle,fill=red!20, node distance=3cm,align=left]
\tikzstyle{teste} = [draw, diamond,fill=green!20, node distance=3cm,align=left]
%    minimum height=2em]
\tikzstyle{phanton} = []   
\tikzstyle{line} = [->,right] %[draw, -latex']
\tikzstyle{retorno} = [loop above]
\tikzstyle{arrow} = [bend left,->]

%  \usepackage[pdftex]{graphicx}  
%  \pdfcompresslevel=1
  % declare the path(s) where your graphic files are
  % \graphicspath{{../pdf/}{../jpeg/}}
  % and their extensions so you won't have to specify these with
  % every instance of \includegraphics
  % \DeclareGraphicsExtensions{.pdf,.jpeg,.png}
\else
  \usepackage[dvips]{graphicx}
  % declare the path(s) where your graphic files are
  % \graphicspath{{../eps/}}
  % and their extensions so you won't have to specify these with
  % every instance of \includegraphics
  % \DeclareGraphicsExtensions{.eps}
\fi
\pdfcompresslevel=9



% *** ALIGNMENT PACKAGES ***
%
\usepackage{array}
% Frank Mittelbach's and David Carlisle's array.sty patches and improves
% the standard LaTeX2e array and tabular environments to provide better
% appearance and additional user controls. As the default LaTeX2e table
% generation code is lacking to the point of almost being broken with
% respect to the quality of the end results, all users are strongly
% advised to use an enhanced (at the very least that provided by array.sty)
% set of table tools. array.sty is already installed on most systems. The
% latest version and documentation can be obtained at:
% http://www.ctan.org/tex-archive/macros/latex/required/tools/


%\usepackage{mdwmath}
%\usepackage{mdwtab}
% Also highly recommended is Mark Wooding's extremely powerful MDW tools,
% especially mdwmath.sty and mdwtab.sty which are used to format equations
% and tables, respectively. The MDWtools set is already installed on most
% LaTeX systems. The lastest version and documentation is available at:
% http://www.ctan.org/tex-archive/macros/latex/contrib/mdwtools/


% IEEEtran contains the IEEEeqnarray family of commands that can be used to
% generate multiline equations as well as matrices, tables, etc., of high
% quality.


\usepackage{eqparbox}
% Also of notable interest is Scott Pakin's eqparbox package for creating
% (automatically sized) equal width boxes - aka "natural width parboxes".
% Available at:
% http://www.ctan.org/tex-archive/macros/latex/contrib/eqparbox/





% *** SUBFIGURE PACKAGES ***
%\ifCLASSOPTIONcompsoc
\usepackage[tight,normalsize,sf,SF]{subfigure}
%\else
%\usepackage[tight,footnotesize]{subfigure}
%\fi
% subfigure.sty was written by Steven Douglas Cochran. This package makes it
% easy to put subfigures in your figures. e.g., "Figure 1a and 1b". For IEEE
% work, it is a good idea to load it with the tight package option to reduce
% the amount of white space around the subfigures. Computer Society papers
% use a larger font and \sffamily font for their captions, hence the
% additional options needed under compsoc mode. subfigure.sty is already
% installed on most LaTeX systems. The latest version and documentation can
% be obtained at:
% http://www.ctan.org/tex-archive/obsolete/macros/latex/contrib/subfigure/
% subfigure.sty has been superceeded by subfig.sty.




% *** FLOAT PACKAGES ***
%
\usepackage{float}
\usepackage{fixltx2e}
% http://www.ctan.org/tex-archive/macros/latex/base/


%\usepackage{stfloats}
% stfloats.sty was written by Sigitas Tolusis. This package gives LaTeX2e
% the ability to do double column floats at the bottom of the page as well
% as the top. (e.g., "\begin{figure*}[!b]" is not normally possible in
% LaTeX2e). It also provides a command:
%\fnbelowfloat
% to enable the placement of footnotes below bottom floats (the standard
% LaTeX2e kernel puts them above bottom floats). This is an invasive package
% which rewrites many portions of the LaTeX2e float routines. It may not work
% with other packages that modify the LaTeX2e float routines. The latest
% version and documentation can be obtained at:
% http://www.ctan.org/tex-archive/macros/latex/contrib/sttools/
% Documentation is contained in the stfloats.sty comments as well as in the
% presfull.pdf file. Do not use the stfloats baselinefloat ability as IEEE
% does not allow \baselineskip to stretch. Authors submitting work to the
% IEEE should note that IEEE rarely uses double column equations and
% that authors should try to avoid such use. Do not be tempted to use the
% cuted.sty or midfloat.sty packages (also by Sigitas Tolusis) as IEEE does
% not format its papers in such ways.




% *** PDF, URL AND HYPERLINK PACKAGES ***
%
\usepackage{url}
% http://www.ctan.org/tex-archive/macros/latex/contrib/misc/
% 
% \url{my_url_here}.




% *** Do not adjust lengths that control margins, column widths, etc. ***
% *** Do not use packages that alter fonts (such as pslatex).         ***
% There should be no need to do such things with IEEEtran.cls V1.6 and later.
% (Unless specifically asked to do so by the journal or conference you plan
% to submit to, of course. )


% correct bad hyphenation here
\hyphenation{op-tical net-works semi-conduc-tor}

\begin{document}
	\title{%\includegraphics[width=18cm]{./fts/cap.png}\\
	\hell\\\ver}
	\author{
%--------------------------------Nomes------------------------------------------
%		\hyperref[luiz]{Luiz Fernando Gomes de Oliveira},~\IEEEmembership{10/46969} 
%		\\\hyperref[panda]{Helbert de Oliveira Coelho Junior},~\IEEEmembership{10/45253}
\ifthenelse{\equal{\journal}{true}}{
\names\\
}{
%\thanks{Revisado em \today.}
\IEEEauthorblockN{\luiz}
\IEEEauthorblockA{Matricula: \luizmatricula\\E-mail: \eluiz}
\and
\IEEEauthorblockN{\canela}
\IEEEauthorblockA{Matricula: \canelamatricula\\E-mail: \ecanela}
\and
\IEEEauthorblockN{\leonardo}
\IEEEauthorblockA{Matricula: \leomatricula\\E-mail: \eleo}
}
\thanks{Revisado em \today.}
	}

%--------------------------O titulo é inserido aqui!----------------------------
%	\selectlanguage{english}  %seleciona o idioma em inglês [para mater os termos abstract e o thanks em inglês]
	\markboth{Universidade de Bras\'ilia - Campus Gama - FGA, \entrega}%
%{	\begin{figure*}
%    	\includegraphics[width=\textwidth]{./fts/unb.jpg}
%	\end{figure*}
%}
	{Shell \MakeLowercase{\textit{et al.}}: \ver}
%-------------------------------------------------------------------------------

	\IEEEcompsoctitleabstractindextext{%
	\begin{abstract}
		% objetivo do experimento
% resumo do procedimento (forma como o experimento foi executado)
% os resultados sucintamente discutidos
% Deve ser feito de preferencia em Inglês
%
%
\selectlanguage{brazil}  %%Desta linha em diante coloque o resumo apenas em português. Nas linhas anteriores, mantenha o texto em inglês.
%
Desenvolver um programa que simule um filtro. O filtro é representado pela média das quatro ultimas amostras e assemelha-se a um filtro passa baixa. O resultado do filtro deverá ser expresso em uma barra de LED's, em displays de 7 segmentos ou em um LCD gráfico.
%


	\end{abstract}
	%%-------------------------------------
	\begin{IEEEkeywords} %palavras chaves
		\keyw %,\LaTeX.
	\end{IEEEkeywords}}
	% make the title area
\onecolumn		%Conferir se pode ser feito em duas colunas

	\maketitle
	\thispagestyle{empty}

	\IEEEdisplaynotcompsoctitleabstractindextext
	\IEEEpeerreviewmaketitle


%\newpage
\pagestyle{headings}
\setcounter{page}{1}
\pagenumbering{roman}

	\selectlanguage{brazil}		%voltando para o português.
	
	\tableofcontents			% Sumario
	\listoffigures
	\listoftables
	\lstlistoflistings

\newpage
\setcounter{page}{1}
\pagenumbering{arabic}
\begin{multicols}{2}    % 2 columns
	\input circuit
\input program


\section{Amostras - laser vermelho}\label{red}

Distancia da fresta ao objeto = \SI{872,0 \pm 1e0}{cm}
Largura do laser no objeto = \SI{11,35 \pm 5e-2}{mm}

\end{multicols}


\begin{enumerate}
	

	\item Verificação da colheira do sinal na fenda dupla.

\input red/dupla_red



	\item Colheita do sinal na fenda dupla - houve uma leve falha no inicio devido à localização do  sensor em frente ao sinal da difração

\input red/simples_red_pissed


	\item Colheira do sinal na fresta simples.

\input red/graph_simples_red


\end{enumerate}

\begin{multicols}{2}    % 2 columns

\section{Amostras - laser verde}\label{green}

\end{multicols}
\begin{enumerate}

\item Fenda simples, laser verde

\input simples

\item Fenda dupla, laser verde

\input dupla
	
\end{enumerate}
\begin{multicols}{2}    % 2 columns

\end{multicols}

%	\ifCLASSOPTIONcaptionsoff
%\newpage
%	\fi

	\bibliographystyle{ieeetr}%{abnt-num}%{ieeetr}%{abnt-alf}
	\bibliography{bibliography}		% expects file "bibliography.bib" 

	%\else
	%\begin{thebibliography}{6} 
	%O Número implica na quantidade MÁXIMA de itens que pode haver de bibliografias.
	%Altere-o se tiver usado mais que duas fontes.

	%\bibitem{IEEEhowto:kopka}
	%H.~Kopka and P.~W. Daly, \emph{A Guide to \LaTeX}, 3rd~ed.\hskip 1em plus
	%  0.5em minus 0.4em\relax Harlow, England: Addison-Wesley, 1999.

	%\bibitem{Wakerly}
	%John F. Wakerly, \emph{Digital Design: Principles and Practices}, 3rd~ed.
	%	\hskip 1em plus 0.5em minus 0.4em\relax Prentice Hall, 1999.

	%\bibitem{zelenovsky}
	%Mendonça A. e Zelenovsky R., \emph{Eletrônica Digital: Curso Prático e Exercícios}, MZ Editora, Brasil, 2004.
	%Contemporary Logic Design - Katz R. H., First Edition, 1993.	

	%\end{thebibliography}
	%\fi

%	\begin{IEEEbiography}[{\includegraphics[width=1in,height=1.25in,clip,keepaspectratio]{./fts/luiz}}]{\luiz}\label{luiz}
%	É, sou eu. Aparecendo aqui só de brinks. Meio que trollando um relatório.
%	\end{IEEEbiography}

\appendices
		\section{Códigos Fontes}\label{src}

%------------------------------------------------------------------------------%		
% Anexar source main.c com link
%\hypertarget{main}{
%\lstinputlisting[language=C]{../main.c} }
%anexar source main.c sem link
%\lstinputlisting[language=C]{../main.c}
%------------------------------------------------------------------------------%		
%\cite{ti_exemplos}
%\tableofcontents %Índice de conteúdos
%\listoftables %Lista de tabelas
%\listoffigures %Lista de figuras
%------------------------------------------------------------------------------%		
\begin{multicols}{2}    % 2 columns
%caption={[ short ] long }

\hypertarget{main_c}{
\lstinputlisting[language=C,caption={[main.c]Código principal do Microcontrolador},label=main_c]{../arduino/main.cpp} }

\vfill\columnbreak
%\newpage
%\end{multicols}
%\vspace{2cm}
%\begin{multicols}{2}    % 2 columns

\hypertarget{main_py}{
\lstinputlisting[language=Python,caption={[uart.py]Código principal de exibição dos dados},label=main_py]{../uart.py} }

\end{multicols}

\onecolumn
\section{Anexos}\label{anexo}

%Inclusão de Figuras
%--------------------Figura logo da UnB----------------------------------------%
%\begin{figure}[h]
%	\centering
%	\includegraphics [scale=1,angle=0,keepaspectratio=true]{./fts/unb}
%	\caption{Logo da UnB}
%	\label{unb}
%\end{figure}
%--------------------Figura logo da UnB----------------------------------------%
%\begin{SCfigure}[1][h] %%Cuidado..ela não se mantem no lugar
%  \centering
%  \includegraphics[width=4cm]{./fts/unb}
%  \caption{Logo da UnB.}
%  \label{unb1}
%\end{SCfigure}


%\begin{multicols}{4}
\tikzstyle{cloud} = [draw, ellipse,fill=blue!20, node distance=3cm,
    minimum height=2em]
\tikzstyle{estado} = [draw, rectangle,fill=red!20, node distance=3cm,align=left]
\tikzstyle{teste} = [draw, diamond,fill=green!20, node distance=3cm,align=left]
%    minimum height=2em]
\tikzstyle{phanton} = []   
\tikzstyle{line} = [->,right] %[draw, -latex']
\tikzstyle{retorno} = [loop above]
\tikzstyle{arrow} = [bend left,->]


\begin{center}
\begin{tikzpicture}[node distance = 2cm]
	\tiny\ttfamily
	%-- Estados
	\node [cloud] (init) at(0,0) {Inicio};
	\node [estado] (E1) at(2,0) {Inicializa\\ UART e ADC};
	\node [estado] (E2) at(4,0) {Captura ADC};
	\node [estado] (E3) at(4,-2) {Escreve ADC na UART};
	%-- Setas
	\path (init) edge [line] (E1);
	\path (E1) edge [line] (E2);
	\path (E2) edge [bend right,->] (E3);
	\path (E3) edge [bend right,->] (E2);
\end{tikzpicture}
\\\hypertarget{diagrama}{Fluxograma do microcontrolador}
\end{center}


\begin{center}
\begin{tikzpicture}[node distance = 2cm]
	\tiny\ttfamily
	%-- Estados
	\node [cloud] (init) at(0,0) {Inicio};
	\node [estado] (E1) at(1,1) {Inicializa\\ serial};
	\node [estado] (E2) at(3,1) {Abre/Cria\\ arquivos};
	\node [teste] (E3) at(5,1) {interrupção\\ do teclado};
	\node [estado] (E4) at(7,1) {Fecha arquivo\\ e serial};
	\node [estado] (E5) at(5,-1) {lê string\\da serial};
	\node [estado] (E6) at(5,-3) {escreve string\\ no arquivo};
	\node [cloud] (fim) at(7,-1) {fim};
	%-- Setas
	\path (init) edge [line] (E1);
	\path (E1) edge [line] (E2);
	\path (E2) edge [line] (E3);
%	\path (E3) edge [line]   (E4); % S
	\draw [line] node at(5.8,1.1) {S} (E3) --  (E4);
%	\path (E3) edge [line] (E5); % N
	\draw [line] node at(5,0.1) {N} (E3) --  (E5);
	\path (E5) edge [line] (E6);
	\path (E6) edge [arrow] (E3);
	\path (E4) edge [line] (fim);
	
\end{tikzpicture}
\\\hypertarget{diagrama}{Fluxograma do computador}
\end{center}


\begin{figure}[H]
\begin{center}\begin{circuitikz} \draw
	(0,0) node[anchor=east]{$V_{ref}$}
	(6,0) node[anchor=west]{$1023-Ad_0$}
	(4.5,-1.5) node[anchor=west]{$Ad_0$}
	(0,0) to [R=$R_1$,*-] (3,3) 
%	(6,0) to [pR=$LDR$, v=$v_1$,*-] (0,0)
	(3,-3) -- node[*-*] {} (3,-1.5)
	(6,0) -- node[*-*] {} (4.5,0)
	(3,-1.5) to [R=$R_4$,*-*] (4.5,0)
	(3,-3) to [R=$R_2$,*-] (0,0)
	(4.5,-1.5) to [R=$R_3$,*-] (3,-3)
	(6,0) to [pR=$LDR$,*-*] (4.5,-1.5)
	(3,3) to [R=$R_5$,*-*] (6,0)
	(3,-3) node[ground] {};
\draw (3,3) -- node[] {} (3,3.5);
\draw (2.5,3.5) --  node[anchor=south] {VCC} (3.5,3.5)  node[anchor=west] {\SI{5}{V}};

; \end{circuitikz} \end{center}
\caption{Idealização do circuito do LDR}
\label{ce1}
\end{figure}

\begin{figure}[H]
\begin{center}
\begin{tikzpicture}[decoration=penciline, decorate,scale=2]
%  \draw[decorate,style=help lines] (-2,-2) grid[step=1cm] (4,4);
	\draw[decorate,thick,pattern=north east lines] (-0.5cm,0.9)   rectangle (5.5,0.5); % barra de baixo
	\draw[decorate,thick,pattern=north east lines] (-0.5cm,4.4)   rectangle (5.5,4.1); % barra de cima
	\draw[decorate,thick,fill=yellow] (0,0) -- (0,5) -- (2,5) -- (2,0) -- (0,0); % primeira chapa
	\draw[decorate,thick,fill=yellow] (3,0) -- (3,5) -- (5,5) -- (5,0) -- (3,0); % segudna chapa
	\draw[decorate,thick,fill=red] (2.4,0) -- (2.4,5) -- (2.6,5) -- (2.6,0) -- (2.4,0); % arame
	\draw[decorate,|-|] (2,-0.3cm) -- (2.3,-0.3cm); % distancia da primeira fenda
	\draw[decorate,|-|] (2.7,-0.3cm) -- (3,-0.3cm); % distancia da segunda fenda
%	\draw[decorate,|-|,dashed] (2.7,-0.3cm) -- (3,-0.3cm); % distancia da segunda fenda
	% notas
	\node[black,scale=0.6] at (1,2.5) {Chapa 1};
	\node[black,scale=0.6] at (4,2.5) {Chapa 2};
	\node[black,scale=0.6] at (2.5,5.2) {Arame};
	\node[black,scale=0.6] at (-1cm,0.7) {Barra de fixação ->};
	\node[black,scale=0.6] at (-1cm,4.2) {Barra de fixação ->};
	\node[black,scale=0.8] at (2.15,-0.5cm) {$\lambda$};
	\node[black,scale=0.8] at (2.85,-0.5cm) {$\lambda$};
\end{tikzpicture}
\end{center}
\caption{Fenda dupla}
\label{fenda_dupla}
\end{figure}

\begin{figure}[H]
\begin{center}
\begin{tikzpicture}[decoration=penciline, decorate,scale=2]
	\draw[thick] (0,0) -- (0,2.4) ; % barra de baixo
	\draw[thick] (0,2.6) -- (0,5) ; % barra de cima
	\draw[thick] (3,0) -- (3,5) ; % barra de choque
	\draw[line] (-1,2.5) -- (-0.2,2.5) ; % barra de choque
	% ondas primarias
	\foreach \A in {0,0.1,0.2,0.3,0.4,0.5} {
		\draw[thick] (-0.8 + \A,2) -- (-0.8 + \A,3) ;
	}
	\node[black,scale=0.8] at (-0.6,1.8) {$\overbrace{\text{Ondas incidentes}}$};
	% ondas após a fenda
	\foreach \A in {0,0.2,0.4,0.6,0.8,1,1.2,1.4} {
		\draw[thick] (0.2+\A,2.0-\A) -- (0.2+\A,2.01-\A) arc (-70:70:0.5+\A) ;
	}
\end{tikzpicture}
\end{center}
\caption{Difração em fenda simples}
\label{diff_simples}
\end{figure}


%--------------------------------
\begin{figure}[H]
\begin{center}
\begin{tikzpicture}[decoration=penciline, decorate,scale=2]
	\draw[thick] (0,0) -- (0,2.2) ; % barra de baixo
	\draw[thick] (0,2.4) -- (0,2.6) ; % barra de baixo
	\draw[thick] (0,2.8) -- (0,5) ; % barra de cima
	\draw[thick] (3,0) -- (3,5) ; % barra de choque
	\draw[line] (-1,2.5) -- (-0.2,2.5) ; % barra de choque
	% ondas primarias
	\foreach \A in {0,0.1,0.2,0.3,0.4,0.5} {
		\draw[thick] (-0.8 + \A,2) -- (-0.8 + \A,3) ;
	}
	\node[black,scale=0.8] at (-0.6,1.8) {$\overbrace{\text{Ondas incidentes}}$};
	% ondas após a fenda
	\foreach \A in {0,0.2,0.4,0.6,0.8,1,1.2,1.4} {
		\draw[thick,color=red] (0.02+\A,2-\A) -- (0.02+\A,2.01-\A) arc (-85:85:0.3+\A) ;
	}
	% ondas após a fenda
	\foreach \A in {0,0.2,0.4,0.6,0.8,1,1.2,1.4} {
		\draw[thick,color=green] (0.02+\A,2.4-\A) -- (0.02+\A,2.41-\A) arc (-85:85:0.3+\A) ;
	}
	% colisão das ondas
	\draw[dashed,->] (0.3,2.2) -- (10:3.5cm) ; % colisão 1
	\draw[dashed,->] (0.31,2.81) -- (52:5.5cm) ; % colisão 2
	\draw[dashed,->] (0.3,2.5) -- (35:4.3) ; % colisão 3
	
	\node[black,scale=0.8] at (4.05,4.35) {Colisão das ondas};
	\node[black,scale=0.8,align=left] at (4.25,2.46) {Colisão das ondas\\ com maior intensidade};
	\node[black,scale=0.8] at (4.1,0.55) {Colisão das ondas};
\end{tikzpicture}
\end{center}
\caption{Difração em fenda dupla}
\label{diff_dupla}
\end{figure}


%\end{multicols}


		
		
\end{document}


